% !TEX program = xelatex
\documentclass{resume}
\usepackage[UTF8,fontset=fandol,scheme=plain]{ctex}
\usepackage[left=17mm,right=17mm,top=13mm,bottom=16mm,footskip=9mm]{geometry}
\setmainfont{texgyrepagella-regular.otf}[
  BoldFont=texgyrepagella-bold.otf,
  ItalicFont=texgyrepagella-italic.otf,
  BoldItalicFont=texgyrepagella-bolditalic.otf]
\usepackage[colorlinks=true,urlcolor=cvAccent,linkcolor=cvAccent]{hyperref}
\hypersetup{pdftitle={Yueci Deng - Curriculum Vitae},pdfauthor={Yueci Deng}}
\AtBeginDocument{\color{cvText}\setlength{\parindent}{0pt}\fontsize{10}{12}\selectfont}

\begin{document}
{\fontsize{24}{27}\selectfont\bfseries Yueci Deng}\hfill
{\small\color{cvMuted}Shenzhen, China}\par
\vspace{3pt}
{\sffamily\fontsize{9.5}{12}\selectfont\color{cvAccent}
Systems, Simulation \& Embodied Intelligence\par}
\vspace{4pt}
{\fontsize{9}{11}\selectfont
\href{mailto:yuecideng@link.cuhk.edu.cn}{yuecideng@link.cuhk.edu.cn}
\enspace\textbar\enspace\href{https://yuecideng.github.io/}{yuecideng.github.io}
\enspace\textbar\enspace\href{https://github.com/yuecideng}{GitHub}
\enspace\textbar\enspace\href{https://scholar.google.com/citations?user=URHu9IUAAAAJ}{Google Scholar}\par}

\begin{rSection}{Education}
\cveducation{The Chinese University of Hong Kong, Shenzhen}{2025 - Present}
{Ph.D. in Computer Science, School of Data Science. Supervisor: Prof. Kui Jia.}
\cveducation{Nanyang Technological University, Singapore}{Jul 2018 - Jun 2019}
{M.Sc. in Signal Processing, Electrical and Electronic Engineering.}
\cveducation{UESTC / University of Glasgow}{Sep 2014 - Jun 2018}
{Joint B.Eng. programme: Electronic Information Engineering (UESTC); Electronics and Electrical Engineering (Glasgow), First Class Honours.}
\end{rSection}

\begin{rSection}{Industry Experience}
\begin{rSubsection}{DexForce Technology Co., Ltd.}{Aug 2022 - Aug 2025}{Head of Core Engine Team}{}
\item Led core engine development for a Sim2Real embodied AI platform, spanning robot simulation and large-scale synthetic data generation.
\item Built training and deployment infrastructure and a unified computer vision and robotics library for product development.
\item Developed physics-based simulation, rendering, and domain-adaptation pipelines for robotic perception and control.
\end{rSubsection}

\begin{rSubsection}{Pudu Technology Inc.}{Dec 2021 - Jul 2022}{Computer Vision Engineer}{}
\item Developed 6D pose estimation, scene reconstruction, and collision avoidance for catering service robots.
\item Automated annotation for synthetic, mixed-reality, and real-world data; researched Sim2Real learning with Tsinghua Shenzhen International Graduate School.
\end{rSubsection}

\begin{rSubsection}{RVBUST Inc.}{Apr 2020 - Dec 2021}{Computer Vision Engineer}{}
\item Developed point cloud registration, recognition, segmentation, and 6D pose estimation for industrial robot vision.
\item Led a cross-functional team developing bin-picking and 3D inspection products.
\end{rSubsection}

\begin{rSubsection}{Appotronics Inc., Research Institute}{Aug 2019 - Apr 2020}{Computer Vision Engineer}{}
\item Developed projector autofocus and calibration algorithms for camera alignment, screen correction, and distortion compensation.
\item Prototyped 3D projection mapping and projected AR systems, and authored related patents.
\end{rSubsection}

\begin{rSubsection}{Panasonic R\&D Centre Singapore}{Sep 2018 - Apr 2019}{Research Engineer}{Internship}
\item Developed object detection, pose estimation, and point cloud processing for robotic picking.
\item Researched deep reinforcement learning for motion planning in simulation. \href{https://github.com/yuecideng/Ur5\_DRL}{[Code]}
\end{rSubsection}

\end{rSection}

\begin{rSection}{Research Grants \& Open Source}
\cvproject{Multimodal Databases for Robot Learning of Human Skills}{Dec 2024 - Dec 2026}
{Principal Investigator \enspace\textbar\enspace RMB 4 million \enspace\textbar\enspace 深圳市具身智能重大专项}
{Lead simulator-based data generation, Real2Sim mapping, and multimodal dataset and skill database development for robot learning.}

\cvproject{\href{https://github.com/DexForce/EmbodiChain}{EmbodiChain}}{}
{Project Lead \enspace\textbar\enspace\href{https://github.com/DexForce/EmbodiChain}{GitHub} \enspace\textbar\enspace\href{https://dexforce.com/embodichain/index.html\#/}{Website}}
{A modular, GPU-accelerated platform for simulation, synthetic data generation, and robot policy training and evaluation, supporting Sim2Real transfer.}

\cvproject{\href{https://github.com/isl-org/Open3D}{Open3D}}{Apr 2022 - Present}
{Open-source Collaborator \enspace\textbar\enspace\href{https://github.com/isl-org/Open3D}{GitHub} \enspace\textbar\enspace\href{https://www.open3d.org/}{Website}}
{Collaborate on tensor-based geometry, 3D reconstruction, and the tensor backend of an open-source 3D data processing library.}
\end{rSection}

\clearpage
{\small\sffamily\color{cvMuted}Yueci Deng\hfill Research \& Technical Profile\par}
% Generated by scripts/sync_cv_publications.py from _pages/about.md.
% Committed with the CV so resume/ remains independently compilable.
\begin{rSection}{Publications and Reports}
{\fontsize{9.5}{11.8}\selectfont
\begin{itemize}[leftmargin=1.8em,labelsep=0.5em,itemsep=5pt,topsep=1pt,parsep=0pt]
\item[{\scriptsize\color{cvMuted}12.}] \href{https://arxiv.org/abs/2610.12333}{\textbf{RiCo: Neural Simulation of Rigid-Body Interactions via Local Contact Reasoning}}\\[-1pt]
{\color{cvMuted}Ruixiang Ouyang, Guanren Qiao, Fansen Meng, \textbf{Yueci Deng}, Ruixing Jin, Kui Jia, Guiliang Liu\textsuperscript{$\dagger$}.} \textit{arXiv preprint,~2026.}
\item[{\scriptsize\color{cvMuted}11.}] \href{https://robosyn-bench.net/material/RoboSynChallenge-report.pdf}{\textbf{RoboSynChallenge: Mastering Real-World Dexterity via Generalizing Synthesized Manipulation Skills}}\\[-1pt]
{\color{cvMuted}RoboSynChallenge Organizing Team.} \textit{NeurIPS Competition Report,~2026.}
\item[{\scriptsize\color{cvMuted}10.}] \href{https://arxiv.org/abs/2602.14965}{\textbf{PAct: Part-Decomposed Single-View Articulated Object Generation}}\\[-1pt]
{\color{cvMuted}Qingming Liu, Xinyue Yao, Shuyuan Zhang, \textbf{Yueci Deng}, Guiliang Liu, Zhen Liu\textsuperscript{$\dagger$}, Kui Jia.} \textit{ACM SIGGRAPH Asia,~2026.}
\item[{\scriptsize\color{cvMuted}9.}] \href{https://arxiv.org/pdf/2606.21525}{\textbf{Backpropagating Through Simulation: Analytic Policy Gradients for Sample and Learning Efficient Differentiable Continuous Control}}\\[-1pt]
{\color{cvMuted}\textbf{Yueci Deng}\textsuperscript{$\dagger$}.} \textit{arXiv preprint,~2026.}
\item[{\scriptsize\color{cvMuted}8.}] \textbf{Focus-Then-Contact: Speeding Up Robotic Contact-Rich Task Learning with Affordance-Guided Real-World Residual Reinforcement Learning}\\[-1pt]
{\color{cvMuted}Guanren Qiao, Ruixiang Ouyang, Sheng Xu, Ruixing Jin, \textbf{Yueci Deng}, Yunxin Tai, Kui Jia, Guiliang Liu\textsuperscript{$\dagger$}.} \textit{ICML,~2026.}
\item[{\scriptsize\color{cvMuted}7.}] \href{https://arxiv.org/abs/2605.11951}{\textbf{From Reaction to Anticipation: Proactive Failure Recovery through Agentic Task Graph for Robotic Manipulation}}\\[-1pt]
{\color{cvMuted}Sheng Xu, Ruixing Jin, Huayi Zhou, Bo Yue, Guanren Qiao, Yunxin Tai, \textbf{Yueci Deng}, Kui Jia, Guiliang Liu\textsuperscript{$\dagger$}.} \textit{RSS,~2026.}
\item[{\scriptsize\color{cvMuted}6.}] \href{https://ieeexplore.ieee.org/document/11495283/}{\textbf{YOTO++: Learning Long-Horizon Closed-Loop Bimanual Manipulation from One-Shot Human Video Demonstrations}}\\[-1pt]
{\color{cvMuted}Huayi Zhou, Ruixiang Wang, Yunxin Tai, \textbf{Yueci Deng}, Guiliang Liu, Kui Jia\textsuperscript{$\dagger$}.} \textit{IEEE TPAMI,~2026.}
\item[{\scriptsize\color{cvMuted}5.}] \href{https://arxiv.org/pdf/2603.17808}{\textbf{EWA: Aligning Video World Models with Executable Robot Actions via Inverse Dynamics Rewards}}\\[-1pt]
{\color{cvMuted}Ruixiang Wang, Qingming Liu, \textbf{Yueci Deng}, Guiliang Liu, Zhen Liu\textsuperscript{$\dagger$}, Kui Jia.} \textit{arXiv preprint,~2026.}
\item[{\scriptsize\color{cvMuted}4.}] \href{https://openreview.net/pdf?id=H4SyKHjd4c}{\textbf{Sim2Real VLA: Zero-Shot Generalization of Synthesized Skills to Realistic Manipulation}}\\[-1pt]
{\color{cvMuted}Runyi Zhao, Sheng Xu, Ruixing Jin, \textbf{Yueci Deng}, Yunxin Tai, Kui Jia, Guiliang Liu\textsuperscript{$\dagger$}.} \textit{ICLR Poster,~2026.}
\item[{\scriptsize\color{cvMuted}3.}] \href{https://www.researchgate.net/profile/Guiliang-Liu-2/publication/396641390\_GS-World\_An\_Efficient\_Engine-driven\_Learning\_Paradigm\_for\_Pursuing\_Embodied\_Intelligence\_using\_World\_Models\_of\_Generative\_Simulation/links/68f33c277d9a4d4e870a9a9d/GS-World-An-Efficient-Engine-driven-Learning-Paradigm-for-Pursuing-Embodied-Intelligence-using-World-Models-of-Generative-Simulation.pdf}{\textbf{GS-World: An Engine-driven Learning Paradigm for Pursuing Embodied Intelligence using World Models of Generative Simulation}}\\[-1pt]
{\color{cvMuted}Guiliang Liu, \textbf{Yueci Deng}, Zhen Liu, Kui Jia\textsuperscript{$\dagger$}.} \textit{Position paper,~2025.}
\item[{\scriptsize\color{cvMuted}2.}] \href{https://openreview.net/pdf?id=AVVXX0erKT}{\textbf{DexScale: Automating Data Scaling for Sim2Real Generalizable Robot Control}}\\[-1pt]
{\color{cvMuted}Guiliang Liu\textsuperscript{*}, \textbf{Yueci Deng}\textsuperscript{*}, Runyi Zhao, Huayi Zhou, Jian Chen, Jietao Chen, Ruiyan Xu, Yunxin Tai, Kui Jia\textsuperscript{$\dagger$}.} \textit{ICML Poster,~2025.}
\item[{\scriptsize\color{cvMuted}1.}] \href{https://arxiv.org/abs/2501.14208}{\textbf{You Only Teach Once: Learn One-Shot Bimanual Robotic Manipulation from Video Demonstrations}}\\[-1pt]
{\color{cvMuted}Huayi Zhou, Ruixiang Wang, Yunxin Tai, \textbf{Yueci Deng}, Guiliang Liu, Kui Jia\textsuperscript{$\dagger$}.} \textit{RSS,~2025.}
\end{itemize}
{\fontsize{8.5}{10}\selectfont\color{cvMuted}* Equal contribution. $\dagger$ Corresponding author.\par}
}
\end{rSection}


\begin{rSection}{Technical Skills}
{\fontsize{9.3}{11.5}\selectfont
\textbf{Programming \& tools}\enspace C / C++ / CUDA, Python, CMake, Git.\\
\textbf{Research \& engineering}\enspace 3D vision, robotics, machine learning, computer graphics, differentiable simulation.\\
\textbf{Delivery}\enspace Applied R\&D, technology productization, large-scale software development.\par}
\end{rSection}

\begin{rSection}{Patents and Applications}
{\fontsize{9.1}{11}\selectfont
\begin{itemize}[itemsep=1pt]
\item \href{https://patents.google.com/patent/CN114155291A/zh}{\textbf{CN114155291A}} --- Box pose recognition method, apparatus, terminal and storage medium.
\item \href{https://patents.google.com/patent/CN114653629B/zh}{\textbf{CN114653629B}} --- Vision-based sorting method, intelligent sorting system and readable storage medium.
\item \href{https://patents.google.com/patent/CN114820781A/zh}{\textbf{CN114820781A}} --- Machine vision-based intelligent handling method, apparatus, system and storage medium.
\item \href{https://patents.google.com/patent/CN217453968U/zh}{\textbf{CN217453968U}} --- Robot, robotic arm and pneumatic gripper.
\item \href{https://patents.google.com/patent/CN121018513A/zh}{\textbf{CN121018513A}} --- Robot URDF model free assembly and constraint management method and system.
\item \href{https://patents.google.com/patent/CN121683799A/zh}{\textbf{CN121683799A}} --- Semantic recognition-based dynamic URDF model construction system and method.
\item \href{https://patents.google.com/patent/CN121934817A/zh}{\textbf{CN121934817A}} --- Robot URDF model construction method, apparatus, electronic device and storage medium.
\item \href{https://patents.google.com/patent/CN122595497A/zh}{\textbf{CN122595497A}} --- Robot model generation method, apparatus, device and medium.
\item \href{https://patents.google.com/patent/CN122574319A/zh}{\textbf{CN122574319A}} --- Grasping pose generation method, apparatus, device and computer-readable storage medium.
\item \href{https://patents.google.com/patent/CN122595647A/zh}{\textbf{CN122595647A}} --- Data processing method, apparatus, computer, storage medium and program product.
\item \href{https://patents.google.com/patent/CN113359381B/zh}{\textbf{CN113359381B}} --- Projection system and projection screen adjustment method.
\item \href{https://patents.google.com/patent/CN113467600B/zh}{\textbf{CN113467600B}} --- Augmented reality-based information display method, system, apparatus and projection device.
\item \href{https://patents.google.com/patent/CN113467601B/zh}{\textbf{CN113467601B}} --- Augmented reality-based information display method, system, apparatus and projection device.
\item \href{https://patents.google.com/patent/CN113497925B/zh}{\textbf{CN113497925B}} --- Projection focusing method and projection focusing device.
\item \href{https://patents.google.com/patent/CN113573035A/zh}{\textbf{CN113573035A}} --- Vision-based adaptive AR-HUD brightness adjustment method.
\item \href{https://patents.google.com/patent/CN113934098B/zh}{\textbf{CN113934098B}} --- Projection device, automatic focusing method and apparatus.
\item \href{https://patents.google.com/patent/CN114286064A/zh}{\textbf{CN114286064A}} --- Real-time focusing method, apparatus, system and computer-readable storage medium.
\end{itemize}
{\fontsize{8.2}{10}\selectfont\color{cvMuted}A: application publication; B: invention grant; U: utility model grant.\par}
}
\end{rSection}
\end{document}
